\documentclass[11pt]{article}
\usepackage[a4paper,margin=27mm]{geometry}
\usepackage[T1]{fontenc}
\usepackage{lmodern,microtype,amsmath,amssymb,amsthm,mathtools,enumitem}
\usepackage[colorlinks=true,linkcolor=blue,citecolor=blue,urlcolor=blue]{hyperref}
\hypersetup{pdftitle={Coherence and Dominance under General Rankings},hypertexnames=false}
\newtheorem{theorem}{Theorem}[section]
\newtheorem{lemma}[theorem]{Lemma}
\newtheorem{proposition}[theorem]{Proposition}
\newtheorem{corollary}[theorem]{Corollary}
\theoremstyle{definition}
\newtheorem{definition}[theorem]{Definition}
\newcommand{\M}{\mathcal M}
\newcommand{\C}{\mathcal C}
\newcommand{\F}{\mathcal F}
\newcommand{\U}{\mathcal U}
\newcommand{\Lf}{\mathfrak L}
\newcommand{\supp}{\operatorname{supp}}
\usepackage{xcolor}
\newcommand{\new}[1]{\textcolor{red}{#1}}
\setlist{itemsep=1pt,topsep=3pt}

\title{Coherence and Dominance under General Rankings}
\author{Jason Konek}
\date{October 2026}

\begin{document}
\maketitle

\begin{abstract}
\end{abstract}

\section{Setting}\label{sec:setting}

\begin{definition}\label{def:rankings}
Let $\Omega$ be a nonempty set of \emph{worlds} and $\M$ a nonempty set of \emph{models}. A \emph{ranking} $L$ consists of a nonempty set $\M_L\subseteq\M$, whose members are the models \emph{evaluable} by $L$, and, for each world $\omega$, a binary relation $\preceq_\omega$ on $\M_L$. We read $m\preceq_\omega m'$ as: $m$ performs at least as well as $m'$ at $\omega$. Models $m$ and $m'$ are \emph{tied at $\omega$} if $m\preceq_\omega m'$ and $m'\preceq_\omega m$. No property of the relations is assumed. $L$ is \emph{transitive} if every $\preceq_\omega$ is transitive. For a fixed ranking $L$, ``model'' means a member of $\M_L$.
\end{definition}

\begin{definition}\label{def:dominance}
Fix a ranking.
\begin{itemize}[leftmargin=*]
\item $c$ \emph{weakly dominates} $m$ if $c\preceq_\omega m$ at every world.
\item $m$ and $c$ are \emph{tied at every world} if each weakly dominates the other.
\item $c$ \emph{dominates} $m$ if $c$ weakly dominates $m$ and they are not tied at every world.
\end{itemize}
A model is \emph{undominated} if no model dominates it. $\U_L$ is the set of undominated models.
\end{definition}

\begin{definition}\label{def:priors}
For a finite nonempty $F\subseteq\Omega$, let $\Delta_F$ be the set of probability mass functions \new{(pmfs)} on $F$, $p=(p_\omega)_{\omega\in F}$, with the Euclidean topology. The support of a pmf $p$ is $\supp p\coloneqq\{\omega:p_\omega>0\}$. Let $\delta_\omega$ be the point mass at $\omega$. When $\Omega$ is finite, $\Delta\coloneqq\Delta_\Omega$.
\end{definition}

\begin{definition}\label{def:coherence-property}
A class $\C\subseteq\M$ of \emph{coherent} models is fixed in advance, independently of the rankings. For a ranking $L$, the \emph{undominated coherent models} of $L$ are
\[
\C_L\coloneqq\{c\in\C\cap\M_L:\ \text{no } m\in\M_L \text{ dominates } c\}=\C\cap\U_L.
\]
%$L$ has the \emph{dominance property} if every model in $\M_L$ is weakly dominated by a member of $\C_L$.
\end{definition}

\subsection{Axioms}

\begin{definition}\label{def:tie-closure}
A ranking $L$ is \emph{closed under ties} if every model tied at every world with a member of $\C_L$ is itself a member of $\C_L$.
\end{definition}

\begin{definition}\label{def:family-assumptions}
Let $\Lf$ be a family of rankings.
\begin{itemize}[leftmargin=*]
\item \emph{Evaluability}: every model in $\M$ is evaluable by some ranking in $\Lf$.
\item \emph{Coherent admissibility}: every coherent model belongs to $\C_L$ for some $L\in\Lf$.
\end{itemize}
\end{definition}

\begin{lemma}\label{lem:ties}
Let $c\in\C_L$ weakly dominate a model $m$.
\begin{enumerate}[label=(\alph*)]
\item If $m$ is undominated, then $m$ and $c$ are tied at every world.
\item If $L$ is closed under ties and $m\notin\C_L$, then $c$ dominates $m$.
\end{enumerate}
\end{lemma}
\begin{proof}
(a) If $m$ did not weakly dominate $c$, then $c$ would dominate $m$. (b) By closure under ties, $m$ is not tied at every world with $c$. Since $c$ weakly dominates $m$, $m$ does not weakly dominate $c$, so $c$ dominates $m$.
\end{proof}

\section{The dominance theorems}\label{sec:general}

\begin{lemma}\label{lem:kkm}
Let $F$ be finite and nonempty, and let $(A_\omega)_{\omega\in F}$ be closed subsets of $\Delta_F$ such that every $p\in\Delta_F$ lies in $A_\omega$ for some $\omega\in\supp p$. Then $\bigcap_{\omega\in F}A_\omega\ne\varnothing$.
\end{lemma}
\begin{proof}
Each $A_\omega$ is nonempty, since by hypothesis we always have 
\[
\delta_\omega\in\bigcup_{\omega^\prime\in\supp \delta_\omega}A_{\omega^\prime}=A_\omega.
\]
Let $d_\omega(p):=\min_{q\in A_\omega}\lVert p-q \rVert$ and $s(p):=\sum_{\omega\in F} d_\omega(p)$. The map $G(p)_\omega:=(\textcolor{red}{p_\omega}+d_\omega(p))/(1+s(p))$ is continuous. Moreover, $G$ maps each $p\in\Delta_F$ to some point in $\Delta_F$ (since each coordinate $G(p)_\omega$ of $G(p)$ is nonnegative and these coordinates sum to 1). Brouwer's fixed-point theorem then guarantees that there is some $p\in\Delta_F$ with $G(p)=p$. This holds just in case $\textcolor{red}{p_\omega}\cdot s(p)=d_\omega(p)$ for all $\omega\in F$. By assumption, there is some $\omega\in\supp p$ with $p\in A_\omega$. Choose some such $\omega$. Then \new{$d_\omega(p)=0$, so} $\textcolor{red}{p_\omega}\cdot s(p)=0$\new{. Since $p_\omega>0$,} $s(p)=0$ and hence $p\in\bigcap_{\omega\in F}A_\omega$.
\end{proof}

This is the lemma of Knaster, Kuratowski and Mazurkiewicz~\cite{KKM}.




\subsection{Finitely many worlds}
\begin{theorem}\label{thm:finite}
Let $\Omega$ be finite and let $L$ be any ranking. Suppose that for each pmf $p\in\Delta$, there is \new{an assigned model} $c_p\in\C_L$. And for each model $m$ and world $\omega$, there is a closed set $A_\omega(m)\subseteq\Delta$, such that:
\begin{enumerate}[label=(\roman*)]
\item $c_p\preceq_\omega m$ for every $p\in A_\omega(m)$;
\item for every model $m$ and every pmf $p\in\Delta$, there is some $\omega\in\supp p$ such that $p\in A_\omega(m)$.
\end{enumerate}
Then:
\begin{enumerate}[label=(\alph*)]
\item for each model $m$, some pmf lies in every $A_\omega(m)$, and the model assigned to any such \new{pmf} weakly dominates $m$;
\item every undominated model is tied at every world with a member of $\C_L$;
\item if $L$ is closed under ties, then $\U_L=\C_L$, and every other model is dominated by a member of $\C_L$.
\end{enumerate}
\end{theorem}
\begin{proof}
Start with (a). Choose $m$. By (ii), the closed sets $A_\omega(m)$ satisfy the hypothesis of Lemma~\ref{lem:kkm}, so some $p^*$ lies in all of them. By (i), $c_{p^*}\preceq_\omega m$ at every world. 

Now proceed to (b) and (c). \new{Let $m$ be any model.} There is some $c\in\C_L$ that weakly dominates $m$, by (a). If $m$ is undominated, Lemma~\ref{lem:ties}(a) ties it to $c$, which gives (b). If \new{$m$ is undominated and} $L$ is closed under ties, then it follows that $m\in\C_L$\new{. As $m$ was an arbitrary undominated model,} $\U_L\subseteq\C_L\subseteq\U_L$. If $m\notin\C_L$, then $c$ dominates $m$ by Lemma~\ref{lem:ties}(b).
\end{proof}

\begin{corollary}\label{cor:closed-choice}
Let $\Omega$ be finite. Suppose that \new{for} each $p\in\Delta$ there is an assigned model $c_p\in\C_L$ \new{such that every model $m$ has some $\omega\in\supp p$ with $c_p\preceq_\omega m$}. Suppose further that for any model $m$ \new{and world $\omega$}, the set $\{p:c_p\preceq_\omega m\}$ is closed. Then the conclusions of Theorem~\ref{thm:finite} hold.
\end{corollary}
\begin{proof}
Take $A_\omega(m)\coloneqq\{p:c_p\preceq_\omega m\}$. Then condition (i) of Theorem~\ref{thm:finite} holds by definition of $A_\omega(m)$, and condition (ii) holds by \new{the first assumption}.
\end{proof}

\subsection{Arbitrary sets of worlds}
%A space is \emph{compact} if every open cover has a finite subcover; no separation property is assumed.

\begin{theorem}\label{thm:local}
Let $\Omega$ be any nonempty set and $L$ any ranking. Suppose:
\begin{enumerate}[label=\arabic*.,leftmargin=*]
\item There is a compact space of parameters $\Theta$. For each $\theta\in\Theta$, there is a selected model $c_\theta\in\C_L$. For each world $\omega$ and model $m$ there is a closed set $K_\omega(m)\subseteq\Theta$ such that $c_\theta\preceq_\omega m$ for every $\theta\in K_\omega(m)$.
\item For every finite set of worlds $G\subseteq\Omega$, there \new{are} a nonempty finite set $F\supseteq G$\new{,} a map $\theta_F:\Delta_F\to\Theta$\new{, and, for each $\omega\in F$ and model $m$, a closed set $A^F_\omega(m)\subseteq\Delta_F$, such that}
\begin{enumerate}[label=(\alph*)]
\item \new{for every $\omega\in F$ and model $m$,} $\theta_F(p)\in K_\omega(m)$ for all $p\in A^F_\omega(m)$;
\item for every model $m$ and every pmf $p\in\Delta_F$, there is some $\omega\in\supp p$ such that $p\in A^F_\omega(m)$.
\end{enumerate}
\end{enumerate}
Then every model is weakly dominated by some $c_\theta$ and conclusions (b) and (c) of Theorem~\ref{thm:finite} hold.
\end{theorem}
\begin{proof}
Choose a model $m$. Let $F$ be any set provided by assumption~2. The closed sets $A^F_\omega(m)$, $\omega\in F$, satisfy the hypothesis of Lemma~\ref{lem:kkm}, by assumption~2(b). So some $p_F$ lies in their intersection. By assumption~2(a), $\theta_F(p_F)\in\bigcap_{\omega\in F}K_\omega(m)$. For any finite number of worlds $\omega_1,\dots,\omega_k$, apply assumption~2 to $G\coloneqq\{\omega_1,\dots,\omega_k\}$ to obtain $F\supseteq G$. Then $\bigcap_iK_{\omega_i}(m)\supseteq\bigcap_{\omega\in F}K_\omega(m)\ni\theta_F(p_F)$. So the closed sets $K_\omega(m)$, $\omega\in\Omega$, have the finite intersection property (every finite subfamily has a non-empty intersection). By compactness some $\theta^*$ lies in all of them, and $c_{\theta^*}\preceq_\omega m$ at every world. Conclusions (b) and (c) follow from Lemma~\ref{lem:ties} as in the proof of Theorem~\ref{thm:finite}.
\end{proof}


\begin{proof}
Choose $m\in\M_L$. For any $\omega\in\Omega$, let $\Theta_\omega=\{\theta\in\Theta:\textcolor{red}{c_\theta\preceq_\omega m}\}$. By (I), each $\Theta_\omega$ is closed. For any $F\in\F$, let
\[
 T_F=\{\theta:\textcolor{red}{c_\theta\preceq_\omega m}\text{ for all }\omega\in F\} = \bigcap_{\omega\in F}\Theta_\omega,
\]
which is closed (as a finite intersection of closed sets). For any $\omega\in F$, let 
\[
A^F_\omega=\sigma_F^{-1}(\Theta_\omega)=\{p\in\Delta_F:\textcolor{red}{c_{\sigma_F(p)}\preceq_\omega m}\}.
\]
These are the sets $A_\omega$ of the proof of Theorem~\ref{thm:kkm}, with $\Delta_F$ in place of $\Delta$ and $c_{\sigma_F(p)}$ in place of $c_p$. They are closed because $\sigma_F$ is continuous, and $\bigcap_{\omega\in F}A^F_\omega=\sigma_F^{-1}(T_F)$. By (S$_\infty$), each $p\in\Delta_F$ lies in $A^F_\omega$ for some $\omega\in\supp p$. So lemma~\ref{lem:kkm} ensures that there is some $p_F\in\bigcap_{\omega\in F}A^F_\omega$. Hence $\sigma_F(p_F)\in T_F$. So $T_F$ is nonempty. For any $F_1,\dots,F_k\in\F$, there is some $G\in\F$ containing their union. In that case, $\varnothing\ne T_G\subseteq T_{F_1}\cap\dots\cap T_{F_k}$. Since $\Theta$ is compact, we must then have $\bigcap_{F\in\F}T_F\neq \varnothing$. (Otherwise the open sets $\Theta\setminus T_F$, $F\in\F$, would cover $\Theta$. By compactness, some finite collection of them, $\Theta\setminus T_{F_1},\dots,\Theta\setminus T_{F_k}$, must cover $\Theta$. But then $T_{F_1}\cap\dots\cap T_{F_k}=\varnothing$, contradicting the claim above.) Choose $\theta_m\in\bigcap_{F\in\F}T_F$ and $\omega\in\Omega$. (S$_\infty$) ensures that there is some $F\in\F$ with $\omega\in F$. Since $\theta_m\in T_F$, \textcolor{red}{$c_{\theta_m}\preceq_\omega m$}. Since $\omega$ was arbitrary, \textcolor{red}{$c_{\theta_m}\preceq m$}, \textit{i.e.}, $c_{\theta_m}$ weakly dominates $m$. The rest of the proof of Theorem~\ref{thm:kkm} applies unchanged, with $c_{\theta_m}$ in place of $c_{p_m}$.
\end{proof}



\section{The dominance argument}\label{sec:conclusion}

\begin{proposition}\label{prop:conclusion}
\new{Let $L$ be a ranking under which every model is weakly dominated by a member of $\C_L$. Then} every undominated model is tied at every world with a member of $\C_L$. If $L$ is also closed under ties, then $\U_L=\C_L$, and every other model is dominated by a member of $\C_L$.
\end{proposition}
\begin{proof}
Apply Lemma~\ref{lem:ties} to each model and a member of $\C_L$ weakly dominating it. By part (a) of the lemma every undominated model is tied with a member of $\C_L$, and so, under closure under ties, lies in $\C_L$; hence $\U_L\subseteq\C_L\subseteq\U_L$. Part (b) of the lemma gives the last claim.
\end{proof}

\begin{proposition}\label{prop:profiles}
Let $L$ be a transitive ranking \new{under which every model is weakly dominated by a member of $\C_L$}. Then $\U_L$ consists of the models tied at every world with some member of $\C_L$, and every other model is dominated by a member of $\C_L$.
\end{proposition}
\begin{proof}
By Proposition~\ref{prop:conclusion}, every undominated model is tied with some $c\in\C_L$. Conversely, let $m$ be tied with $c\in\C_L$, and suppose $m'$ dominates $m$. Then $m'\preceq_\omega m\preceq_\omega c$ at every world, so $m'$ weakly dominates $c$ by transitivity. If $c$ weakly dominated $m'$, then $m\preceq_\omega c\preceq_\omega m'$ would make $m$ weakly dominate $m'$, which is false. So $m'$ dominates $c$, which is impossible. Hence $m$ is undominated. Finally, if $m$ is dominated, let $c\in\C_L$ weakly dominate it; $m$ is not tied with $c$, so $c$ dominates $m$.
\end{proof}

\begin{definition}\label{def:permissibility}
Let $\Lf$ be a family of rankings. A model is \emph{irrational} if every ranking in $\Lf$ that evaluates it has an undominated model that dominates it. A model is \emph{rationally permissible} if it is undominated under some ranking in $\Lf$.
\end{definition}

\begin{proposition}\label{prop:permissible}
Suppose that \new{under each ranking $L\in\Lf$, every model in $\M_L$ is weakly dominated by a member of $\C_L$}. Then for each $L\in\Lf$:
\begin{enumerate}[label=(\alph*)]
\item every undominated model is tied at every world with a member of $\C_L$;
\item if $L$ is closed under ties, then $\U_L=\C_L$, and every other model in $\M_L$ is dominated by a member of $\C_L$;
\item if $L$ is transitive, then $\U_L$ consists of the models tied at every world with a member of $\C_L$, and every other model in $\M_L$ is dominated by a member of $\C_L$.
\end{enumerate}
If every ranking in $\Lf$ is closed under ties or transitive, then a model is undominated under no ranking iff, under every ranking that evaluates it, some undominated model dominates it. Coherent admissibility alone, without \new{this assumption}, makes every coherent model \new{rationally} permissible. If, in addition, every ranking in $\Lf$ is closed under ties, the \new{rationally} permissible models are exactly the coherent ones.
\end{proposition}
\begin{proof}
Parts (a) and (b) are Proposition~\ref{prop:conclusion}, and (c) is Proposition~\ref{prop:profiles}. If $m\notin\bigcup_L\U_L$, then for each $L$ that evaluates $m$, by (b) or (c), some member of $\C_L\subseteq\U_L$ dominates $m$. The converse is immediate from the definitions. Coherent admissibility says $\C\subseteq\bigcup_L\C_L\subseteq\bigcup_L\U_L$. If every ranking is closed under ties, (b) gives $\bigcup_L\U_L=\bigcup_L\C_L$, which is $\C$ by coherent admissibility, since every $\C_L\subseteq\C$.
\end{proof}

\begin{thebibliography}{9}
\bibitem{KKM} B.~Knaster, C.~Kuratowski and S.~Mazurkiewicz. Ein Beweis des Fixpunktsatzes f\"ur $n$-dimensionale Simplexe. \emph{Fund. Math.} 14 (1929) 132--137.
\end{thebibliography}
\end{document}
